% Generated draft: requires mathematical and source review.
\documentclass[11pt]{article}
\usepackage[margin=25mm]{geometry}
\usepackage{fontspec}
\setmainfont{Latin Modern Roman}
\usepackage{amsmath,amssymb,mathtools,unicode-math}
\setmathfont{Latin Modern Math}
\usepackage{enumitem,needspace,xurl,hyperref}
\setlength{\parindent}{0pt}
\setlength{\parskip}{5pt}
\setlength{\emergencystretch}{3em}
\newcommand{\N}{\mathbb{N}}
\newcommand{\Z}{\mathbb{Z}}
\newcommand{\Q}{\mathbb{Q}}
\newcommand{\R}{\mathbb{R}}
\newcommand{\C}{\mathbb{C}}
\newcommand{\F}{\mathbb{F}}
\newcommand{\A}{\mathbb{A}}
\newcommand{\PP}{\mathbb{P}}
\newcommand{\E}{\mathbb{E}}
\newcommand{\eps}{\varepsilon}
\newcommand{\dd}{\,\mathrm{d}}
\begin{document}
\section*{UnsolvedMath ID: 800677}
\textbf{Source problem code:} catalogue entry
\subsection{Definitions and mathematical statement}
Let \(\mathbb P^2=\mathbb P^2(\C)\). A leaf of a singular holomorphic foliation \(\mathcal F\) means a maximal connected integral Riemann surface in the regular locus of \(\mathcal F\), equipped with its intrinsic complex structure. Write
\[
\mathbb D:=\{\zeta\in\C:|\zeta|<1\}.
\]
Let \(\mathscr B_d\) denote the space of holomorphic foliations on \(\mathbb P^2\) of geometric degree \(d\). Here geometric degree has its standard meaning: the number of tangencies, counted with multiplicity, with a generic projective line.

Fix an affine coordinate chart \(\C^2\subset\mathbb P^2\), with coordinates \((z,w)\), and put \(L_\infty:=\mathbb P^2\setminus\C^2\). As in the supplied source, let \(\mathscr A_d\) denote the space of foliations represented in this chart by a polynomial vector field
\[
Z=P(z,w)\frac{\partial}{\partial z}
  +Q(z,w)\frac{\partial}{\partial w},
\qquad P,Q\in\C[z,w],\qquad
\max\{\deg P,\deg Q\}=d.
\]
The source explicitly distinguishes this polynomial degree from geometric degree and states that the polynomial degree is not invariant under automorphisms of \(\mathbb P^2\). No identification of \(\mathscr A_d\) with \(\mathscr B_d\) is assumed. The source does not specify normalization of polynomial representatives, including treatment of common polynomial factors of \(P\) and \(Q\).

\textbf{Question (Anosov, as presented by Dinh and Sibony, Problem 4).} For a generic foliation \(\mathcal F\in\mathscr A_d\), are generic leaves of \(\mathcal F\) biholomorphic to \(\mathbb D\)? Ask the same question, separately, for a generic foliation \(\mathcal F\in\mathscr B_d\). For a leaf \(L\), the requested conclusion is the existence of a biholomorphic map
\[
f:L\longrightarrow\mathbb D.
\]
The quantifier pattern is generic foliations followed, for each such foliation, by generic leaves; the conclusion concerns the leaf itself.

\textbf{Supplied background and unresolved qualifications.} The record states that generically in \(\mathscr B_d\) foliations have no compact leaf, and generically in \(\mathscr A_d\) the line \(L_\infty\) is invariant. These are retained as source assertions, without independent verification. The record does not define either meaning of \emph{generic}, specify a topology or measure for these quantifiers, or state the intended range of \(d\). Consequently, this draft preserves the question but cannot replace its genericity language with fully specified quantifiers or impose an unstated restriction on \(d\).
\subsection{Short English statement}
For each of the two supplied families of foliations on \(\mathbb P^2(\C)\), ask whether a generic foliation has generic leaves biholomorphic to the unit disc. Polynomial degree and geometric degree are distinct conventions. The intended degree range, the two meanings of genericity, and normalization of polynomial representatives require clarification.
\subsection{Sources}
\begin{itemize}
\item Local UnsolvedMath registry, permanent ID \(800677\), exact source problem code \texttt{catalogue entry}; supplied complete source record in the AMR Open Problem Lists collection. The registry attributes the presentation to Dinh and Sibony, with proposed year 2018, and labels the problem \emph{open}. This is a supplied status label, not an independently verified literature assessment.
\item Dinh and Sibony, \emph{Some Open Problems on Holomorphic Foliation Theory} (2018), Problem 4, \url{https://arxiv.org/abs/1807.05640}. The supplied record gives Anosov as attribution/context. It records extraction by \texttt{arxiv\_latex} and access on 2026-10-16; these are supplied provenance details, not a claim of access or verification during this formalization.
\item Supplied upstream revision: \texttt{d371ee917ac9647fbb06cd6fc529deec6f3dda82}. The registry records creation and update on 2026-10-16.
\end{itemize}
\end{document}
